\cotefile{46}
\scanpage{65}
% archive sheet number: 64
% L1-5
$a\in\unsure{\mathscr{B}|\mathcal{B}}(X)$, soit $a^\circ$ la classe ``opposée'' (on remplace une algèbre par l'algèbre opposée), alors $aa^\circ$ est trivial. En effet, on a canoniquement $A\otimes A^\circ\simeq L(A,A)$ : D'où la définition du \uline{groupe de Brauer}, noté $\operatorname{Br}(X)$.

% L6-9
On \unsure{note|voit} qu'on a
\[
\delta(ab)=\unsure{I_m|\iota_m}\delta(a)+\unsure{I_n|\iota_n}\delta(b)
\]
\margin{compte tenu des hom
$\begin{gathered}\mathbb{Z}/n\mathbb{Z}\\\mathbb{Z}/m\mathbb{Z}\end{gathered}\longrightarrow\mathbb{Z}/mn\mathbb{Z}$}
d'où résulte qu'on a un hom de \deleted{\unsure{groupes|groupe}} \inserted{monoïdes}
\[
\mathscr{B}(X)\longrightarrow H^2(X,\mathbb{Q}/\mathbb{Z})
\]
% L10-14
Cet homomorphisme définit, par composition
$\unsure{B|\mathscr{B}}(X)\longrightarrow H^2(X,\underline{\mathbb{C}}^*)$, \deleted{\unsure{d'où}} et \deleted{\unsure{est trivial}} \inserted{son noyau est} \deleted{\unsure{sur les}\ \illegible{2}} \inserted{formé des} éléments ``triviaux''. \deleted{\unsure{on}} D'où un hom naturel \uuline{injectif}
\[
\boxed{\operatorname{Br}(X)\xrightarrow{\unsure{\varphi|\psi}}H^2(X,\underline{\mathbb{C}}^*)}
\]
% L15-16
dont l'image est contenue dans $\operatorname{Im}H^2(X,\mathbb{Q}/\mathbb{Z})$, donc dans le sous-groupe de Torsion.

% L17-20
\uline{Question 1} Tout élément de torsion de $H^2(X,\underline{\mathbb{C}}^*)$ provient-il d'un élément de $\operatorname{Br}(X)$ ?

\unsure{Cela|Ce} résulterait de \unsure{l'affirmative|l'affirmation} à la question suivante plus précise

% L21-22
\uline{Question 2} Tout élément de $H^2(X,\mathbb{Z}/n\mathbb{Z})$ est-il de la forme $\delta a$, avec $a\in H^1(X,\unsure{\operatorname{GP}|\operatorname{Gp}}(n,\ ))$ ???

% The source leaves the second GP argument blank.
\scanpage{66}
% archive sheet number: 65
% L1-4
[Je voudrais d'abord résoudre cette question du point de vue \uline{topologique}. \unsure{Le cas crucial du|Le cas crucial est le} cas de la classe canonique de $H^2(2,\mathbb{Z}/n\mathbb{Z};\mathbb{Z}/n\mathbb{Z})$, est-elle de la forme $\delta(a)$, avec $a\in H^1(\unsure{2},\mathbb{Z}/n\mathbb{Z},\underline{\operatorname{GP}(n-1)})$ ?
\margin{\unsure{faisceau} des \unsure{sections|unités} continues}
% L5-7
Sinon, on aurait le \unsure{corollaire|conséquence} de 2 suivant : $H^2(X,\mathbb{Q}/\mathbb{Z})$, tout élément est-il l'obstruction d'un élément de $H^1(X,\underline{\operatorname{GP}(n-1)})$ pour $n$ convenable ??

% L8-13
Puis on peut $=$ \unsure{se ramener|se demander} à voir si tout élément de torsion de $H^3(X,\mathbb{Z})$ $\deleted{\unsure{\simeq H^2(X}}$ \deleted{\unsure{$\delta$}} [i.e. tout élément de $\partial(H^2(X,\mathbb{Q}/\mathbb{Z}))$] est l'obstruction d'un élément d'un $H^1(X,\operatorname{GP}(n-1))$, avec $n$ convenable [N.B. $H^3(X,\mathbb{Z})\simeq H^2(X,\underset{\sim}{\mathbb{C}}^*)$] --- c'est la question 1 \uline{topologique} --- ].

% L14-16
On peut introduire le revêtement universel $\bar X$ de $X$, de groupe de Galois $\pi$, et on a alors
\[
0\unsure{{}-{ }|\longrightarrow}H^2(\pi,\mathbb{Z}/n\mathbb{Z})\unsure{{}-{ }|\longrightarrow}H^2(X,\mathbb{Z}/n\mathbb{Z})\longrightarrow H^2(\bar X,\mathbb{Z}/n\mathbb{Z})^\pi\longrightarrow\cdots
\]
% L17-23
et on peut essayer de montrer que du moins les éléments de $H^2(X,\mathbb{Z}/n\mathbb{Z})$ qui proviennent de $H^2(\pi,\mathbb{Z}/n\mathbb{Z})$ sont dans $\operatorname{Im}H^1(X,\operatorname{GP}(n-1))$ ; \deleted{\unsure{c. à d.}} il suffit pour ceci de prouver que $H^1(\pi,\operatorname{GP}(n-1))\to H^2(\pi,\mathbb{Z}/n\mathbb{Z})$ est surjectif --- \deleted{\unsure{ce qui}\ \illegible{4}\ \inserted{\unsure{pourquoi}}\ \unsure{quand} $X=K(1,\pi)$}. Du moins, est-il vrai que $\forall$ élément de $H^2(\pi,\mathbb{Q}/\mathbb{Z})$ est dans $\operatorname{Im}H^1(\pi,\operatorname{GP}(n-1))$ pour $n$ \uline{convenable} ?
\scanpage{67}
% archive sheet number: 66
% L1-3
Et si ceci aussi est faux, est-il \unsure{donc} vrai \unsure{du|au} moins que tout élément de $H^2(\pi,\mathbb{C}^*)$ est dans $\operatorname{Im}\deleted{\unsure{\operatorname{GP}(\mathbb{C})|\operatorname{GP}(n)}}H^1(\pi,\operatorname{GP}(n-1))$ pour $n$ convenable ??

% L4: three short dashes
\noindent $-\quad-\quad-$
\margin{]}
\scanpage{68}
% archive sheet number: 67
% L1-5, title at upper right
``\uline{Restriction des scalaires}'' dans les groupes algébriques et leurs \unsure{espaces} principaux \unsure{homogènes}.
\scanpage{69}
% archive sheet number: 68
% L1-3
$k$ corps, $k'$ extension finie

$S=\operatorname{Spec}k$, $S'=\operatorname{Spec}k'$

\margin{N.B. les \unsure{mêmes|énoncés} sont vrais : voir dans la \unsure{généralisation} ---}

\unsure{a)} $X'$ préschéma sur $k'$, $X_1=\prod_{S'/S}X'/S'$.

% L4
\[
\boxed{H^0(k,X_1)\simeq H^0(k',X)}\tag*{(i)}
\]
% L5-6
(ii) Si $X'$ est un schéma en groupes, soit $G'$, \uline{\unsure{lisse} sur $k$}, on a
\[
\boxed{H^1(k,G_1)\simeq H^1(k',G')}
\]
(Classes de fibrés principaux homogènes ---)

% L7-8
(iii) Si $G'$ est \deleted{\unsure{un}} un schéma en groupe \uline{comm.}, soit $\unsure{M'|M}$, on a
\[
\boxed{H^i(k,M_1)\simeq H^i(k',M')}\qquad\text{pour }\forall i
\]
(Cohomologie galoisienne\unsure{...|.})

% L9
(i) Est vrai par définition (\unsure{on s'en fout du corps},

% L10-18
(ii) Soit $P'$ fibré principal homogène sous $G'$, alors $P_1=\prod_{S'/S}P'$ est un schéma sur $k$, sur lequel $G'$ opère. On voit que $P_1$ est formellement fibré pr. hom. sous $G_1$, d'autre \deleted{\unsure{part il est}\ \illegible{2}} part utilisant la \unsure{lissité|simplicité} de $G'/k'$, on voit que $P_1\ne\varnothing$, \deleted{\unsure{que c'est un fibré principal}} \inserted{\unsure{donc est un fibré principal}} homogène.
% slanted
\margin{\unsure{Contre-exemple} si $G'$ est comm. \unsure{radiciel} ?? Si $G'$ est comm. \unsure{ou propre|non propre|non lisse}/\unsure{$k'$} ??}
Le foncteur $P'\mathrel{\unsure{\rightsquigarrow|\longmapsto}}P_1$ est \uline{pleinement fidèle}.
\margin{(i.e. $\unsure{H^1|H_1|H P_1}$ est \unsure{bijectif})}
Pour montrer que c'est une équivalence de catégories, un argument de descente ramène au cas où $P_1$ est trivial ---

% L19-22
(iii) Soit $T/S$ \deleted{fini plat} quelconque, alors on a
\[
C^*(T/S,M_1)\simeq C^*(T'/S',M')
\]
d'où
\[
H^{\unsure{\alpha|q}}(T/S,M_1)\simeq H^\alpha(T'/S',M'),
\]
d'où, si $T$ parcourt certains $S$-préschémas
\[
\varinjlim_T H^\alpha(T/S,M_1)\simeq\varinjlim_T H^\alpha(T'/S',M'),
\]
% L23-26
Or, si $T$ parcourt les rev. étales (\unsure{ou p.ex.} les \unsure{préschémas} étales séparés quasi-finis surjectifs sur $S$) les $T'/S'$ sont \uline{cofinaux} dans la \inserted{catégorie} \unsure{des revêtements|des recouvrements} rel. à $S'$.

D'où la conclusion

\margin{TSVP}
\scanpage{70}
% archive sheet number: 69
% L1-2
Supposons que $k'/k$ \unsure{soit|est} \uline{radiciel}, prenons $\deleted{\unsure{A'}}M'=G_{m\,k'}$, on a alors une suite exacte
% L3-4
\diagram{\begin{tikzcd}[ampersand replacement=\&,column sep=large]
\unsure{e|0}\arrow[r,no head]\&G_m\arrow[r]\&(G_m)_1\arrow[r]\arrow[d,equal]\&U\arrow[r]\&e\\
\&\&\underline{\operatorname{Hom}}_S(S',G_m)\&\&
\end{tikzcd}}
% L5-9
où $U$ est unipotent comm. \unsure{lisse|simple} sur $k$ [en car. 2, on peut trouver $U$ de dim 1, \deleted{\unsure{p.ex.}} si $[k':k]=2$]. La suite exacte de Coh. galoisienne sur $k$ donne alors, compte tenu de $H^{\unsure{i|1}}(k,(G_m)_1)\simeq H^i(k',G_m)$, l'isom
\[
H^1(k,U)\simeq\operatorname{Ker}(H^2(k,G_m)\unsure{{}-{ }|\longrightarrow}H^2(k',G_m))
\]

% L10-13
Prenons p.ex. $k=k_0[[t]]$, $k'=k_0[[t^{1/p}]]$, où $k_0$ est un corps fini de car. $p$, on trouve par le \unsure{corps de classes local|corps de classes locale}
\[
H^1(k,U)\simeq\mathbb{Z}/p\mathbb{Z}\qquad\ne0
\]
% slanted
\margin{N.B. On voit facilement que \unsure{ici on a} $pU=0$ (vrai \unsure{chaque fois} \unsure{pour|que} $k'^p\subset k$)}
% L14-15
[bien que, pour $p=2$, $U$ soit une forme de $G_a$ ; bien sur, cette forme n'est pas isom. à $G_a$].

% L16-19
Notons aussi que \uline{l'ext. de} $U$ \uline{par} $G_m$ [qui \unsure{splitte} sur $\bar k$] \uline{n'est pas triviale}, car autrement on aurait
\[
\underset{\substack{\shortparallel\\0}}{H^1(k,(G_m)_1)}\simeq\underset{\substack{\shortparallel\\0}}{H^1(k,G_m)}\times H^1(k,U)
\]
% L20-21
d'où $H^1(k,U)=0$. Donc ``\uline{le \unsure{r.|v.} \inserted{\uline{unipotent}} de} $(G_m)_1$ \uline{n'est pas défini sur} $k$''
\scanpage{71}
% archive sheet number: 70
% L1-3
Regardons maintenant l'\uline{hom. norme}, et soit $V$ son noyau
\[
e\longrightarrow V\longrightarrow(G_m)_1\mathbin{-}G_m\longrightarrow e
\]
% L4-8
\uline{Mais}, $V$ \uline{n'est pas lisse sur $k$}, \deleted{\unsure{le}\ \illegible{1}\ \unsure{groupe}} \deleted{Mais [dans le cas où \unsure{donc}} pas plus que $V_{\mathrm{réd}}$ [car autrement ce serait le radical \inserted{\uline{unipotent}} de $G_m$].
\margin{dans le cas précédent envisagé où $H^1(k,U)\ne0$}
D'ailleurs, dans le cas où $[k':k]=p$, on voit facilement que $V$ est réduit. D'ailleurs, dans le cas où $k'^p\subset k$, on voit facilement que $V$ \uline{est aussi le noyau de la puissance $p$-ième dans} $G_{m\,1}$, donc on a un isom
% slanted
\margin{\uuline{\unsure{préciser}} \unsure{prouver par descente} ---}
\[
V\simeq\prod_{S'/S}(\unsure{\boldsymbol{\mu}|\mu}_p)_{S'}
\]
% L13-17
\deleted{Dans tous cas}, ce qui prouve qu'on a en fait une suite exacte
\[
0\mathbin{-}\boldsymbol{\mu}_p\mathbin{-}V\longrightarrow U\longrightarrow0
\]
(extension non triviale ---).

% L18-22
Notons maintenant que si $G'$ est un groupe alg. \inserted{\unsure{connexe simple|connexe réduit}} sur $k'$ \deleted{\unsure{un groupe}} ayant un centre isom. à $(\boldsymbol{\mu}_p)_{k'}$
\margin{[p.ex. $\operatorname{SL}(p)\deleted{\unsure{(k')}}_{k'}$]}
alors le centre de $G_1$ est
\[
Z_1=V
\]
donc $(\bar Z_1)_{\mathrm{réd}}$ \uline{n'est pas définissable sur} $k$
\scanpage{72}
% archive sheet number: 71
% L1-2, title in upper-right corner
\uline{Petits fourbis \unsure{chomologiques|cohomologiques}}
\scanpage{73}
% archive sheet number: 72
% margin note written sideways
\margin{Petits fourbis \unsure{chomologiques|cohomologiques}}
% L1-9
\uline{Proposition} Soit $k$ un corps commutatif, $\unsure{L/k}$ une extension galoisienne de $k$ de groupe de Galois $\Gamma$, $\underline A$ une algèbre avec 1 de dimension finie sur $k$, $\underline A^L$ l'algèbre sur $L$ obtenue par extension des scalaires à $L$, $A^L$ est considéré comme muni du groupe d'automorphismes d'anneau $\Gamma$. Soit $G=\underline A^{L*}$ le groupe des éléments inversibles de $\underline A^L$, $G$ \deleted{\unsure{$\underline A^L$|est}} un $\Gamma$-groupe.

$H^1(\Gamma,G)$ est réduit à l'élément neutre, ainsi que $H^1(\Gamma,G^\psi)$ si $\psi$ est un hom croisé de $\Gamma$ dans $\deleted{\operatorname{Int}G}\operatorname{Aut}\underline A$.

% L10
\uline{Corollaire} $H^1(\Gamma,\operatorname{Gl}(n,L))$ est réduit à l'élément neutre.

% L11-12
\uline{Démonstration} $\Gamma$ opère dans $G^\psi=G$

$\sigma(\gamma)g=\psi(\gamma).g^\gamma$. Soit $\unsure{\varphi|\phi}\in Z^1(\Gamma,G^\psi)$, \deleted{\unsure{posons}} faisons
% L13
\[
\deleted{\unsure{a}}\quad x=\sum_{\gamma\in\Gamma}\bigl(\deleted{\unsure{\sigma(\gamma).A}}\bigr)\varphi(\gamma)\sigma(\gamma).A\qquad(A\in\underline A^L)
\]
% L14, cancelled line
\deleted{les \unsure{$\varphi(\gamma)$}\ \unsure{les|et les}\ $\sigma(\gamma)$ étant \unsure{en fait}\ \illegible{1}}
\deleted{\unsure{$\exists$}}\deleted{\illegible{1}}
% L15-18
\[
\begin{aligned}
\sigma(\gamma)x&=\sum_{\gamma'\in\Gamma}\bigl(\deleted{\unsure{\sigma(\gamma\gamma')A}}\bigr)\sigma(\gamma)\varphi(\gamma')\sigma(\gamma\gamma').A\\
&\deleted{=\sum_{\gamma'\in\Gamma}(\sigma(\gamma')A)\sigma(\gamma)\varphi(\gamma^{-1}\gamma')}\\
&=\sum_{\gamma'\in\Gamma}\sigma(\gamma)\varphi(\gamma^{-1}\gamma')\sigma(\gamma')A\\
&=\sigma(\unsure{\gamma})\varphi(\gamma^{-1})\sum_{\gamma'\in\Gamma}\varphi(\gamma')\sigma(\gamma')A=\sigma(\gamma)\varphi(\gamma^{-1})x
\end{aligned}
\]
% L19
\[
x=\varphi(\gamma^{-1})\sigma(\gamma)^{-1}x\qquad x=\varphi(\gamma)\sigma(\gamma)x
\]
% L20-21
$\deleted{\unsure{\varphi(\gamma)=}}$ D'où, \uline{si $x$ inversible}, la formule
\[
\varphi(\gamma)=x\bigl(\sigma(\gamma)x\bigr)^{-1}
\]
% L22-23
Il reste à prouver donc que l'on peut trouver $A$ de telle façon que $\sum\varphi(\gamma)\sigma(\gamma).A$ soit
\scanpage{74}
% archive sheet number: 73
% L1-6
inversible. Or $\sigma(\gamma)A=\psi(\gamma).A^\gamma$, où $\psi(\gamma)$ est une transformation linéaire de $\underline A^L$, si c'était non inversible, on aurait \deleted{\unsure{on}}

($P$ étant la fonction \inserted{polynôme} $P(A)=\deleted{\unsure{\det(\rho_A)}}$ \deleted{\unsure{de la}\ \illegible{1}} \deleted{\unsure{fonction de}} déterminant de la transformation $B\unsure{\longrightarrow|\text{---}}AB$ resp. $B\unsure{\longrightarrow|\text{---}}BA$)
\[
P\Bigl(\sum_\gamma \unsure{\varphi|\phi}(\gamma)\bigl(\psi(\gamma).A^\gamma\bigr)\Bigr)=0
\]
% L7-9
Montrons qu'il est impossible que ceci soit \unsure{bien} \deleted{$=0$ \unsure{pour}\ $P(\sum_\gamma$} pour $\unsure{\forall|\exists} A$. Considérons la fonction sur $(\underline A^L)^\Gamma$ définie par
\[
\deleted{\unsure{\overline P((A_\gamma))=P(\sum_\gamma\varphi(\gamma))}}
\]
% L10
\[
\overline P((A_\gamma))=P\Bigl(\sum_\gamma\varphi(\gamma)\bigl(\psi(\gamma)A_\gamma^{\deleted{\unsure{\gamma}}}\bigr)\Bigr)
\]
% L11-13
Elle est nulle si \inserted{on fait la substitution} $A_\gamma=A^\gamma$ $(A\in\underline A^L)$ d'où on conclut qu'elle est identiquement nulle (du moins si \uline{le corps $k$ est infini})
% slanted
\margin{\unsure{indépendance} algébrique des \unsure{automorphismes} d'un corps!!!}
% L14-16
On aura donc, en faisant $A_\gamma=0$ si $\gamma\unsure{=|\ne}1$, $A_1=1\in\underline A^L$ : $P(1)\deleted{\unsure{\varphi(1)}}=0$, ce qui est absurde.

% L17-20
\unsure{Reste} : \unsure{Le pb du cas|Le cas du corps} où $k$ est \uline{fini}. Dans ce cas, \deleted{$G$ est \unsure{un groupe}\ \unsure{fini}} \unsure{considérons le groupe} $G$ \unsure{algébrique sur $k$ et comme} \deleted{\unsure{fini}}\deleted{\unsure{un $\Gamma$-groupe}}, \deleted{\unsure{soit $u$ un générateur de $\Gamma$}.} \uline{\unsure{Pour le prouver|par le groupe}}, et prouvons $H^1(\Gamma,G^{\deleted{\unsure{\psi}}\psi})=e$ si $\psi$ est
% slanted; insertion continued in the left margin
\inserted{\unsure{comme ci-dessus} un hom. croisé de $\Gamma$ dans $\operatorname{Aut}_{\unsure{L|k}}G^{\unsure{L}}$}
% slanted
\margin{$\Gamma$ est cyclique, $u\lambda=\lambda^q$ $(q=\unsure{[k]|\lvert k\rvert})$ un générateur.}
$\varphi$ est un \unsure{hom croisé} \unsure{connu si} $\varphi(u)=a$
% L21-25
\[
\begin{aligned}
\varphi(u^0)&=1\\
\varphi(u)&=a\\
\varphi(u^2)&=a\,\sigma(u)a\\
\varphi(u^3)&=a\,\sigma(u)a\quad\sigma(u^2)a\\
\overline\varphi(u^{n-1})&=a\,\sigma(u)a\ \cdots\ \sigma(u)^{n-2}a
\end{aligned}
\]
\scanpage{75}
% archive sheet number: 74
% L1-6
\uline{Proposition} Soit $G$ un groupe algébrique \uline{connexe} sur un corps \uline{fini} $k$, $L$ une extension finie de $k$, de groupe de Galois $\Gamma$, $G^L$ le groupe obtenu à partir de $G$ par extension des scalaires à $L$ (il est encore connexe sur $L$) $\Gamma$ opère dans $G^L$ comme groupe d'automorphismes. \deleted{Soit} \deleted{\unsure{$\psi$ un}\ \illegible{1}} Alors $H^1(\Gamma,G^L)=\{e\}$. Plus précisément,
% slanted; alongside a vertical emphasis stroke, not a cancellation
\margin{\unsure{6|b}\quad\unsure{pas vrai!}}
% L7-11
soit $\operatorname{Aut}G^L$ le groupe des automorphismes du groupe algébrique $G^L$ \uline{sur $L$}, et soit $\psi$ un \deleted{\unsure{groupe}}\inserted{hom} croisé de $\Gamma$ dans $\operatorname{Aut}G^L$, soit $G^{L\psi}$ le groupe $G^L$ \deleted{\unsure{muni}} où $\Gamma$ opère par $\sigma(\gamma).g=\psi(\gamma).g^\gamma$ ; alors on a encore
\[
H^1(\Gamma,G^{L\psi})=\{e\}
\]

% L12-15
\uline{Démonstration} Soit $q=\text{nb d'éléments de }k$, et $n=[L:k]$.
$\Gamma$ est le groupe cyclique engendré par l'automorphisme $\unsure{f|u}\lambda=\lambda^q$, on désigne par $u$ l'automorphisme correspondant de $G^L$. Soit $\Omega$ une clôture algébrique de $L$, on désigne par $\widetilde u$ l'automorphisme $\deleted{\unsure{\widetilde f\lambda=\lambda^q\text{ de }\Omega}}$, \deleted{\unsure{et par $\widetilde u$ aussi}} \deleted{$\unsure{x\longrightarrow x^{\widetilde u}}$} \deleted{l'automorphisme} \deleted{$x\longrightarrow x^q$} de $G^\Omega$ \deleted{\unsure{défini ainsi}} défini \inserted{par l'automorphisme $\lambda\longrightarrow\lambda^q$ de $\Omega$} (noter que $u^n=1$, mais $\widetilde u^n\ne1$ ; l'ensemble des invariants de $\widetilde u^n$ dans $\Omega$ --- resp$^t$ $G^\Omega$ --- est $L$ --- resp. $G^L$ ---).
% L16-20
On désigne par $\widetilde\sigma$ l'automorphisme $x\longrightarrow\psi(f)x^{\widetilde u}$ de $G^\Omega$, induisant un automorphisme $\sigma$ de $G$, qui n'est autre que $\sigma(f)$. On a, puisque $\psi(\unsure{\gamma|f})$ est un \deleted{$\unsure{\psi\,\psi^{\widetilde u}\psi^{\widetilde u^2}\cdots\psi^{\widetilde u}}$} hom croisé de $\Gamma$ dans $\operatorname{Aut}G^L$
\scanpage{76}
% archive sheet number: 75
% L1
\[
\psi^{(u)}\psi^{(u^2)}\cdots\psi^{(u^{n-1})}=1
\]
% L2-5
(où $\psi^{(u^k)}=u^k\psi u^{k-1}$), d'où on conclut, en désignant par $\widetilde\psi$ le prolongement de $\psi$ à $G^\Omega$, et notant que les $\psi^{(u^k)}$ sont encore des $\unsure{k|L}$-automorphismes \inserted{du gr. algébrique} $G^L$ par la \deleted{\unsure{restriction}}\inserted{restriction} des automorphismes $\widetilde\psi^{(\widetilde u^k)}$ de $G^\Omega$ :
% L6
\[
\widetilde\psi^{\widetilde u}\widetilde\psi^{\widetilde u^2}\cdots\widetilde\psi^{\widetilde u^{n-1}}=1 \tag*{(*)}
\]
% L7-8
Soit $\unsure{\varphi|\phi}\in Z^1(\Gamma,G^{L\psi})$, donné par $\varphi(u)=a\in G^{L\unsure{\psi}}$, avec
\[
a\,a^{\unsure{\sigma|v}}a^{\sigma^2}\cdots a^{\sigma^{n-1}}=1 \tag*{(**)}
\]
% L9-11
On cherche $b\in G^L$ tel que $a=b(b^\sigma)^{-1}$ ; on va voir qu'il suffit de \unsure{y|à} trouver $b\in G^\Omega$ tel que $a=b(b^{\widetilde\sigma})^{-1}$ : on va voir qu'on aura automatiquement $b\in G^L$.
% L12-15
Pour ceci, on \deleted{\unsure{note que}} \deleted{\unsure{part de}\ $a=b(b^{\widetilde\sigma})^{-1}$} : \unsure{dans} $(**)$
\[
\deleted{\unsure{a^{\sigma^k}|a^{u^k}}=a^{\widetilde\sigma^k}}
\]
\[
\deleted{\widetilde\sigma^k=(\widetilde\psi\widetilde u)^k=\widetilde\psi\widetilde\psi^{\widetilde u}\widetilde\psi^{\widetilde u^2}\cdots\widetilde\psi^{\widetilde u^{\unsure{k|k-1}}}\widetilde u^k}
\]
\deleted{(récurrence sur $k$)}
\deleted{D'où}
\[
\deleted{a^{\sigma^k}=a^{\widetilde\sigma^k}=(\widetilde\psi\cdots\widetilde\psi^{\widetilde u^{k-1}})\,b\,b^\sigma}
\]
\[
\deleted{\unsure{b^{\widetilde\sigma-1}|b^{\widetilde\sigma^{-1}}}}.
\]
% L16-18
en notant que
\[
a^{\sigma^k}=a^{\widetilde\sigma^k}=b^{\widetilde\sigma^k}(b^{\widetilde\sigma^{k+1}})^{-1}.
\]
d'où on tire
\[
b(b^{\widetilde\sigma^n})^{-1}=1\qquad\text{soit}\qquad b^{\widetilde\sigma^n}=b
\]
% L19-21
Or on a
\[
\widetilde\sigma^k=(\widetilde\psi\widetilde u)^k=\widetilde\psi\widetilde\psi^{\widetilde u}\cdots\widetilde\psi^{\widetilde u^{k-1}}\widetilde u^k
\]
(récurrence sur $k$) d'où
\[
\widetilde\sigma^n=(\widetilde\psi\widetilde\psi^{\widetilde u}\cdots\widetilde\psi^{\widetilde u^{n-1}})\widetilde u^n=\widetilde u^n
\]
% L22-24
en vertu de $(*)$, on a donc $b^{\unsure{\widetilde u^n}}=b$, i.e. $b\in G^L$.

On est maintenant ramené au théorème géométrique suivant :
\scanpage{77}
% archive sheet number: 76
% L1-5
\uline{Lemme!} Soit $G$ un groupe algébrique connexe sur le corps algébriquement clos $\Omega$ de car. $p\ne0$, \deleted{soit $q$ une puissance de $p$} \inserted{``défini'' sur le corps $\unsure{\mathbb F_q|F_q}$ à $q$ éléments, soit} $\psi$ un automorphisme \inserted{du groupe algébrique} \deleted{de} $G$, considérons \deleted{l'endomorphisme du groupe $G$\unsure{,|.} défini par} \deleted{\unsure{abstrait}} l'application de $G$ dans $G$ donnée par
% L6
\[
\deleted{\unsure{\theta|\varphi}}\qquad x\longrightarrow x(\psi x^{(q)})^{-1}
\]
% L7-9
où $x^{(q)}$ est le transformé de $x$ par l'automorphisme qui correspond à $\lambda\longrightarrow\lambda^{(q)}$ dans $\Omega$ (dont le corps des inv. est $\mathbb F_q$) cette application est alors \uline{surjective}.

% L10-14
\deleted{On considère l'application $x\longrightarrow$}
Considérons la différentielle de cette application à l'élément neutre : par un calcul purement formel, on trouve que c'est l'identité (car la différentielle de $x\unsure{\text{---}|\longrightarrow}\psi x^{(q)}$ est la \deleted{\unsure{transformée}}\inserted{composée} de $d\psi(e)$ avec la différentielle de $x^{(q)}$, laquelle est nulle)
% L15-18
Donc l'application envisagée a pour image un \uline{voisinage} de $e$. C'est d'ailleurs un hom. croisé (quand $G$ opère sur lui $=$ par automorphismes intérieurs par $\psi x^{(q)}$ :
\[
\unsure{\Theta|\theta}(xy)=\Theta(x)\operatorname{int}(\psi x^{(q)})\Theta(y).
\]
% L19-21
Si $G$ est abélien, c'est que $\Theta(G)$ est donc un sous-groupe, qui est \inserted{ouvert donc} identique à $G$ puisque $G$ est connexe.

Dans le cas général, on \deleted{\unsure{a}}
\scanpage{78}
% archive sheet number: 77
% L1-6
Soit $G$ un groupe où opère $\mathbb Z$, \inserted{${\unsure{\Pi|T}}^1=T$.} Pour $\forall$ entier $n$, soit $G_n$ l'ens des éléments de $G$ tels que $\Pi^n x=x$.
On suppose que $G=\bigcup G_n$ (les $G_n$ forment une suite croissante). Alors $G_n$ est un $\mathbb Z_n$-groupe, et on peut considérer sa cohomologie \inserted{$H^i$} (\unsure{si $i>1$, supposons que} $G$ commutatif), alors on a
% L7-11
\[
H^0(\mathbb Z_n,G_n)=G_1\qquad\text{ens. des \unsure{él$^{ts}$|él$^t$} de $G$ invariants par $\mathbb Z$}
\]
\[
\begin{aligned}
H^1(\mathbb Z_n,G_n)&=\text{ens. des éléments $\alpha$ de $G_n$ tels que}\\
&\quad\alpha\,T\alpha\,T^2\alpha\cdots T^{n-1}\alpha=1,\quad\text{mod la relation d'équivalence}\\
&\quad\text{définie par }\beta\simeq\alpha\text{ si et seulement si}\\
&\quad\text{il existe }\unsure{u|x}\in G_n\text{ tel que }\beta=u^{-1}\alpha\,Tu
\end{aligned}
\]
% L12-14
D'ailleurs, supposons que $n\mid m$, alors $\mathbb Z_n$ est un quotient de $\mathbb Z_m$, \deleted{\unsure{et on a}} par un sous-groupe $N_{n,m}$, et on a
\[
\boxed{G_n=G_m^{N_{nm}}}.
\]
On a donc fonctoriellement une suite exacte
% L15
\diagram{
\begin{tikzcd}[ampersand replacement=\&]
e\arrow[r]\& H^1(\mathbb Z_n,G_n)\arrow[r,"\text{inj}"]\&H^1(\mathbb Z_m,G_m)\arrow[r]\&H^1(\underbrace{N_{nm}}_{\substack{\Vert\\\mathbb Z_{m/n}}},G_m)^{\unsure{\mathbb Z_m|\mathbb Z_n}}
\end{tikzcd}}
% L16-20
[Si \deleted{$\unsure{G_n}$} $G$ est abélien, on a aussi $H^i(\mathbb Z_n,G_n)$, et on a des hom
\[
H^i(\mathbb Z_n,G_n)\longrightarrow H^i(\mathbb Z_m,G_m)
\]
d'ailleurs, on a à partir de $i=1$, une \uline{périodicité} de période 2 dans les $H^i(\mathbb Z_n,G_n)$, \deleted{\unsure{pour les}} \unsure{les $i>1$ n'y gagnent pas grand chose}].

% L21-23
\uline{Lemme!} Soit $\alpha\in G$ tel que $\alpha\,T\alpha\,T^2\alpha\cdots T^{n-1}\alpha=1$, alors $\alpha\in G_n$.
Supposons $\alpha,\beta\in G$ satisfaisant aux conditions précédentes, et $\beta=u^{-1}\alpha\,Tu$ $(u\in G)$, alors $u\in G_n$.
% L24-27
En effet, \unsure{on aura}, appliquant $T$ \deleted{$\unsure{\cdot}$}
\[
T\alpha\,T^2\alpha\cdots T^{n-1}\alpha\,T^n\alpha=1
\]
i.e. $T^n\alpha$ est un \uline{inverse} \deleted{\illegible{1}} de $T\alpha\,T^2\alpha\cdots T^{n-1}\alpha$ ; comme il en est de $=$ de $\alpha$, on a $T^n\alpha=\alpha$.
De $\unsure{\equiv|=}$, écrivons $\beta\,T\beta\cdots T^{n-1}\beta=1$, en y remplaçant $\beta$ par $u^{-1}\alpha\,Tu$, on trouve
\scanpage{79}
% archive sheet number: 78
% L1-2, bibliographic notes crossed together by a wavy stroke
\deleted{F. \unsure{Peterson|Petterson} Functional Cohomology operations. Some non-embedding Problems.}

\deleted{Moore The double suspension and $p$-primary components of homotopy gps of spheres.}
% L3-5; cancellation strokes through paired factors
\[
u^{-1}\alpha\,\deleted{Tu\,Tu^{-1}}\,T\alpha\,\deleted{T^2u\,T^2u^{-1}}\,T^2\alpha\,\deleted{T^3u}\cdots\deleted{T^{n-1}u^{-1}}\,T^{n-1}\alpha\,T^nu=1
\]
i.e.
\[
u^{-1}(\alpha\,T\alpha\cdots T^{n-1}\alpha)T^nu=1
\]
d'où, comme $(\quad)=1$, $u^{-1}T^nu=1$ soit $T^nu=u$.
% L6-9
\uline{Lemme 2} Si un élément de $H^1(\mathbb Z_n,G_n)$ est représenté par un $\alpha\in G_n$ tel que $\alpha\,T\alpha\cdots T^{n-1}\alpha=1$, alors son image dans $H^1(\mathbb Z_m,G_m)$ (si $n\mid m$) est de $=$ représentée par $\alpha\in G_n\subset G_m$.

\unsure{Trivial|Triviale}
% L10
Posons
\[
{\unsure{\mathscr H|\mathcal H|\mathbb H}}^{\unsure{1}}(T,G)=\underset{\substack{\unsure{n|\longrightarrow}\\\text{suivant l'ordre}\\\text{filtrant arithmétique}}}{\lim}H^1(\mathbb Z_n,G_n).
\]
% L11-14
On voit donc que $H^1(T,G)$ est l'ens des classes d'objets $\alpha\in G$ tels qu'il existe un $n$ avec $\alpha\,T\alpha\cdots T^{n-1}\alpha=1$, \deleted{\unsure{et}} en écrivant $\beta\sim\alpha$ s'il existe $u\in G$ tel que $\beta=u^{-1}\alpha\,Tu$.
% L15-17
\uline{Lemme 3} Si \unsure{Pour} \deleted{\unsure{tout}\ \unsure{avec|groupe}\ $G$} l'application $x\unsure{\longrightarrow|\text{---}}x^{-1}Tx$ de $G$ dans $G$ est surjective, on a $H^1(T,G)=\{e\}$.

Facile sur lemmes 1 et 2
% L18-23
\uline{Lemme 4} Soit
\diagram{\begin{tikzcd}[ampersand replacement=\&]
e\arrow[r,dash]\&G'\arrow[r,dash]\&G\arrow[r,dash]\&G''\arrow[r,dash]\&e
\end{tikzcd}}
une suite exacte de $\mathbb Z$-groupes ($T$ opère dans $G$ et conserve $G'$).
Supposons que pour $\forall$ automorphisme \inserted{intérieur $\sigma$} de $G$, \deleted{\unsure{on a}} \deleted{\unsure{pour chaque hom intérieur de $G$ l'application|pour chaque automorphisme intérieur de $G$ l'application}} \inserted{l'image de $G$ par l'application} $x\unsure{\longrightarrow|\text{---}}x^{-1}\sigma Tx$ \deleted{de $G'$ dans $G'$ soit surjective} \inserted{contienne $G'$}. Alors
\[
H^1(T,G)\unsure{\longrightarrow|\text{---}}H^1(T,G'')\quad\text{est \uline{injectif}}
\]
% L24-25
Soient en effet $\alpha,\beta\in G$, tels que $\alpha\,T\alpha\cdots T^{n-1}\alpha=1$ et de $=$ $\beta\,T\beta\cdots T^{n-1}\beta=1$, et
\scanpage{80}
% archive sheet number: 79
% L1-7
supposons qu'il existe $u\in G$ tel que $\beta\equiv u^{-1}\alpha\,Tu$ $(G')$.
On veut prouver que $\beta\sim\alpha$. Remplaçant $\alpha$ par $u^{-1}\alpha\,Tu$, on peut supposer déjà que $\beta\equiv\alpha$ $(G')$, i.e. $\beta=t\alpha$ $(t\in G')$.
On cherche $v\in G'$ tel que $\beta=v^{-1}\alpha\,Tv$ i.e. $t\alpha=v^{-1}\alpha\,Tv$ i.e. $t=v^{-1}\alpha\,Tv\alpha^{-1}$ ou encore, ($\sigma$ désigne l'aut. de $G'$ induit par l'aut. intérieur défini par $\alpha$) $t=v^{-1}\sigma T(v)$. Un tel $v$ existe par hypothèse.
% L8-12
On cherche maintenant des conditions moyennant lesquelles $H^1(T,G)\unsure{\longrightarrow|\text{---}}H^1(T,G'')$ est surjectif. Soit donc $\alpha\in G$ tel que $\alpha\,T\alpha\cdots T^{n-1}\alpha\equiv1$ $(G')$, on veut montrer qu'on peut trouver $t\in G'$ tel que $\beta=t\alpha$ satisfasse à $\beta\,T\beta\cdots T^{n-1}\beta\unsure{\equiv|=}1$. Or le premier mb est
% L13-14; the right extremities touch the scan edge
\[
\begin{aligned}
&(t\alpha\,T(t)T(\alpha)\cdots T^{n-1}(t)T^{n-1}\unsure{(\alpha)|\alpha})\\
&=t\underbrace{(\alpha\,T(t)\alpha^{-1})}\underbrace{(\alpha\,T\alpha\,T^2(t)T(\alpha)^{-1}\alpha^{-1})}\cdots\underbrace{(\alpha\,T\alpha\cdots T^{n-2}\alpha\,T^{n-1}(t)T^{n-2}\alpha^{-1}\cdots T\alpha^{-1}\alpha^{-1})}\underbrace{\alpha\,T\alpha\cdots T^{n-1}\alpha}
\end{aligned}
\]
% L15-18
Le dernier facteur est 1, et on trouve
\[
t\sigma_1T(\unsure{t|\theta})\sigma_2T^2(t)\cdots\sigma_{n-1}T^{n-1}(t)
\]
où $\sigma_i$ est l'automorphisme intérieur défini par $\alpha\,T\alpha\cdots T^{i-1}\alpha$. D'où le \deleted{\unsure{lemme}}
% L19-24
\uline{Lemme 5} Soit
\diagram{\begin{tikzcd}[ampersand replacement=\&]
e\arrow[r,dash]\&\unsure{G'}\arrow[r,dash]\&G\arrow[r,dash]\&G''\arrow[r,dash]\&e
\end{tikzcd}}
une suite exacte de $\mathbb Z$-groupes, on suppose que pour \unsure{$\forall$e|toute} suite d'automorphismes $\sigma_1\cdots\sigma_{n-1}$ de $G'$, induits par des automorphismes intérieurs de $G$, l'application
\[
x\longrightarrow x\sigma_1Tx\cdots\sigma_{n-1}T^{n-1}x
\]
de $G$ dans $G$ est \uline{surjective}.
Alors
\[
H^1(\unsure{X|x},G)\unsure{\longrightarrow|\text{---}}H^1(X,G'')
\]
